% Volume I: The Epistemology of Suspicion
% Part II. Events, Records, and Reconstructions
% Chapter 7: Provenance and the Chain of Transformation
% Spine: proof-spines/ch007-spine.md

\chapter{Provenance and the Chain of Transformation}
\label{ch:ch007}


Once the event is separated from its record, provenance must be modeled as a transformation history rather than a badge of origin. A record's provenance is represented as a directed acyclic graph whose nodes are transformations and whose edges carry custody relations. Confidence therefore attaches not simply to a final artifact, but to the integrity of the path by which it was observed, stored, transmitted, edited, and retrieved.

For a path \(p=(e_1,\ldots,e_k)\), one provisional integrity measure is
\[
I(p)=\prod_{j=1}^{k}r(e_j),
\]
where \(r(e_j)\in[0,1]\) measures the reliability of transformation \(e_j\). The chapter uses that formalism only as a starting point, since multiplicative trust can conceal correlated failure and common-source dependence. The real lesson is that evidential confidence must be distributed across the chain of transformation, not projected backward from the last visible artifact.
